\documentclass[11pt,a4paper]{article}

% Template (template/labs.sty, template/tikz-redes.sty)
\makeatletter
\def\input@path{{../../../template/}{template/}}
\makeatother
\usepackage{labs}          % opcoes: gabarito (versao do professor), nofira
\usepackage[brazil]{babel}
\usepackage{amsmath}

% Versao do professor: descomente a linha abaixo (ou use \usepackage[gabarito]{labs})
% \gabaritotrue

\labdisciplina{Interconexão de Redes de Computadores}
\labnumero{Revisão de roteamento}
\labtitulo{Rotas estáticas e sumarização}
\labmodulo{Revisão}
% \labencontro{12}
% \labduracao{60 min}
\labferramentas{Papel e caneta}

% Linha em branco para preenchimento nas tabelas de rotas
\newcommand{\lv}{\rule{0pt}{2.0em} & & & \\ \hline}

\begin{document}
\cabecalholab

\begin{objetivos}
  \item Montar a tabela de rotas de cada roteador de uma topologia com três equipamentos, distinguindo rotas conectadas de rotas estáticas.
  \item Decidir, a partir dos blocos de endereços já alocados, quais redes podem ser anunciadas por uma única rota resumida e quais precisam de rotas específicas.
  \item Verificar o alinhamento de blocos CIDR por meio da representação binária do último octeto.
\end{objetivos}

\begin{prerequisitos}
  \item Endereçamento IPv4 e notação CIDR.
  \item Entrega direta e indireta, tabela de roteamento e \termo{longest prefix match} (encontro anterior).
\end{prerequisitos}

\newpage
\section{Contexto}

O CEEI está expandindo a infraestrutura de rede de laboratórios distribuídos em dois blocos do campus: o REENGE e o CAA. A equipe de TI do CEEI reservou o bloco \codigo{10.20.4.0/22} para essa expansão e já definiu o endereçamento de cada laboratório. Você, como responsável técnico pela implantação, recebe os blocos prontos e precisa planejar a conectividade entre os laboratórios e o núcleo de rede do CEEI, decidindo como anunciar essas redes da forma mais enxuta possível.

Cada prédio tem um roteador na entrada: \termo{R1} no REENGE, \termo{R2} no CAA. Os dois se conectam a \termo{R3}, o roteador do CEEI, que concentra o tráfego dos laboratórios antes de seguir para o restante do campus. Você vai montar as rotas para o cenário atual e, mais adiante neste roteiro, revisá-las quando novos laboratórios entrarem em operação nos dois prédios.

\section{Cenário atual}

\begin{center}
\begin{tikzpicture}
  % Roteadores
  \node[roteador, label={[rotulo]right:{R3 (CEEI)}}]  (r3) at (0,2.8) {};
  \node[roteador, label={[rotulo]left:{R1 (REENGE)}}] (r1) at (-4.4,0) {};
  \node[roteador, label={[rotulo]right:{R2 (CAA)}}]   (r2) at (4.4,0) {};

  % Laboratorios do REENGE
  \host{lcc}{(-5.4,-2.6)}{LCC\\10.20.4.0/26}
  \host{labarc}{(-3.4,-2.6)}{LABARC\\10.20.4.64/26}

  % Laboratorios do CAA
  \host{laeg}{(3.4,-2.6)}{LAEG\\10.20.5.0/27}
  \host{mat}{(5.4,-2.6)}{Lab. Materiais\\10.20.5.64/28}

  % Enlaces com o CEEI
  \draw[tronco] (r1) -- node[rotuloenlace]{10.20.0.0/30} (r3);
  \draw[tronco] (r2) -- node[rotuloenlace]{10.20.0.4/30} (r3);

  % Enlace de contingencia (fora do escopo deste roteiro)
  \draw[wan] (r1) -- node[rotuloenlace,below]{contingência 10.20.0.8/30} (r2);

  % Enlaces de laboratorio
  \draw[enlace] (r1) -- (lcc);
  \draw[enlace] (r1) -- (labarc);
  \draw[enlace] (r2) -- (laeg);
  \draw[enlace] (r2) -- (mat);

  % Dominios
  \begin{scope}[on background layer]
    \node[dominio, fit={(r1)(lcc)(labarc)($(lcc.west)+(-0.5,0)$)($(labarc.east)+(0.5,0)$)($(labarc.south)+(0,-1.1)$)}] (dreenge) {};
    \node[dominio, fit={(r2)(laeg)(mat)($(laeg.west)+(-0.5,0)$)($(mat.east)+(0.5,0)$)($(mat.south)+(0,-1.1)$)}] (dcaa) {};
    \node[dominio, fit={(r3)($(r3.north)+(0,0.5)$)($(r3.west)+(-1.4,0)$)($(r3.east)+(2.2,0)$)}] (dceei) {};
  \end{scope}
  \node[dominiotitulo] at (dreenge.north west) {REENGE \quad 10.20.4.0/24};
  \node[dominiotitulo] at (dcaa.north west)    {CAA \quad 10.20.5.0/24};
  \node[dominiotitulo] at (dceei.north west)   {CEEI (agregação)};
\end{tikzpicture}
\end{center}

O enlace de contingência entre R1 e R2 existe fisicamente, mas \destaque{não deve ser usado neste roteiro}: considere apenas os enlaces com o CEEI. A rota de contingência será tratada no roteiro prático.

\subsection{Endereçamento}

Os blocos abaixo já estão definidos pela equipe de TI. O gateway de cada laboratório é o primeiro endereço utilizável do bloco, configurado na VLAN correspondente do roteador do prédio.

\begin{center}
\renewcommand{\arraystretch}{1.25}
\begin{tabular}{@{}llll@{}}
\toprule
\cabtab{Prédio} & \cabtab{Laboratório} & \cabtab{Bloco} & \cabtab{Gateway} \\
\midrule
REENGE & LCC (Ciência da Computação)        & 10.20.4.0/26   & 10.20.4.1   \\
REENGE & LABARC (Arquitetura de Computadores) & 10.20.4.64/26  & 10.20.4.65  \\
\midrule
CAA    & LAEG (Apoio ao Estudante da Graduação) & 10.20.5.0/27   & 10.20.5.1   \\
CAA    & Lab. de Materiais                  & 10.20.5.64/28  & 10.20.5.65  \\
\bottomrule
\end{tabular}
\end{center}

\begin{center}
\renewcommand{\arraystretch}{1.25}
\begin{tabular}{@{}lllp{5.4cm}@{}}
\toprule
\cabtab{Enlace} & \cabtab{Rede} & \cabtab{Lado A} & \cabtab{Lado B} \\
\midrule
R1 (REENGE) a R3 (CEEI) & 10.20.0.0/30 & R1: 10.20.0.1 & R3: 10.20.0.2 \\
R2 (CAA) a R3 (CEEI)    & 10.20.0.4/30 & R2: 10.20.0.5 & R3: 10.20.0.6 \\
\bottomrule
\end{tabular}
\end{center}


Dicas para a próxima etapa;

\begin{conceito}[Rota estática e resumo]

Duas redes podem ser \termo{resumidas} (agregadas) em uma única rota quando o bloco resultante cobre exatamente as duas, sem incluir endereços que não pertencem a nenhuma delas. Para dois blocos de mesmo tamanho \codigo{/n}, isso acontece quando eles são vizinhos e o primeiro está alinhado ao início de um bloco \codigo{/(n-1)}. O teste prático é olhar o último octeto em binário: o resumo é o prefixo comum às duas redes.
\end{conceito}

\begin{conceito}[Rota conectada]
Uma \termo{rota conectada} aparece na tabela automaticamente quando uma interface (do switch, uma VLAN) recebe endereço IP: o roteador sabe entregar diretamente naquela rede. Uma \termo{rota estática} é escrita à mão e diz "para chegar nesta rede, entregue ao próximo salto (\termo{next-hop}) tal". Quando um pacote casa com mais de uma rota, vence a de prefixo mais longo (\termo{longest prefix match}).

% Duas redes podem ser \termo{resumidas} (agregadas) em uma única rota quando o bloco resultante cobre exatamente as duas, sem incluir endereços que não pertencem a nenhuma delas. Para dois blocos de mesmo tamanho \codigo{/n}, isso acontece quando eles são vizinhos e o primeiro está alinhado ao início de um bloco \codigo{/(n-1)}. O teste prático é olhar o último octeto em binário: o resumo é o prefixo comum às duas redes.
\end{conceito}

% \begin{cuidado}[Restrição deste roteiro]
% Não use rota padrão (\codigo{0.0.0.0/0}) em nenhum roteador. Toda rede alcançável deve ter uma entrada explícita na tabela, seja conectada, seja estática. O objetivo é exercitar a decisão de resumir ou não em cada roteador, e a rota padrão esconderia essa decisão.
% \end{cuidado}

\subsection{Questão}
Preencha a tabela de rotas de cada roteador. Em \textbf{Tipo}, escreva \emph{conectada} ou \emph{estática}; para rotas estáticas, informe o next-hop. Use quantas linhas precisar.

\begin{tarefa}{Tabela de rotas de R1 (REENGE)}
\renewcommand{\arraystretch}{1.0}
\begin{tabular}{|p{3.6cm}|p{2.6cm}|p{2.6cm}|p{5.2cm}|}
\hline
\cabtab{Rede de destino} & \cabtab{Next-hop} & \cabtab{Tipo} & \cabtab{Observação} \\ \hline
\lv \lv \lv \lv \lv \lv \lv
\end{tabular}
\end{tarefa}

\begin{gabarito}
\begin{tabular}{@{}lllp{5.4cm}@{}}
\toprule
\cabtab{Destino} & \cabtab{Next-hop} & \cabtab{Tipo} & \cabtab{Observação} \\
\midrule
10.20.4.0/26   & \cinza{(VLAN LCC)}    & conectada & \\
10.20.4.64/26  & \cinza{(VLAN LABARC)} & conectada & \\
10.20.0.0/30   & \cinza{(enlace CEEI)} & conectada & \\
10.20.5.0/27   & 10.20.0.2 & estática & LAEG, via R3 \\
10.20.5.64/28  & 10.20.0.2 & estática & Lab. Materiais, via R3 \\
10.20.0.4/30   & 10.20.0.2 & estática & opcional: enlace R2 a R3, só para testes com ping \\
\bottomrule
\end{tabular}

As duas redes do CAA não admitem resumo exato (ver Pergunta 1), então R1 precisa de uma rota estática para cada uma. Aceite também a resposta que resume o CAA em \codigo{10.20.5.0/24} (o bloco inteiro do prédio), desde que o aluno explique que a rota cobre endereços não utilizados.
\end{gabarito}

\begin{tarefa}{Tabela de rotas de R2 (CAA)}
\renewcommand{\arraystretch}{1.0}
\begin{tabular}{|p{3.6cm}|p{2.6cm}|p{2.6cm}|p{5.2cm}|}
\hline
\cabtab{Rede de destino} & \cabtab{Next-hop} & \cabtab{Tipo} & \cabtab{Observação} \\ \hline
\lv \lv \lv \lv \lv \lv \lv
\end{tabular}
\end{tarefa}

\begin{gabarito}
\begin{tabular}{@{}lllp{5.4cm}@{}}
\toprule
\cabtab{Destino} & \cabtab{Next-hop} & \cabtab{Tipo} & \cabtab{Observação} \\
\midrule
10.20.5.0/27   & \cinza{(VLAN LAEG)}      & conectada & \\
10.20.5.64/28  & \cinza{(VLAN Materiais)}  & conectada & \\
10.20.0.4/30   & \cinza{(enlace CEEI)}     & conectada & \\
10.20.4.0/25   & 10.20.0.6 & estática & resumo de LCC + LABARC, via R3 \\
10.20.0.0/30   & 10.20.0.6 & estática & opcional: enlace R1 a R3 \\
\bottomrule
\end{tabular}

Uma única rota resume todo o REENGE. Aceite também duas rotas específicas (\codigo{/26} cada), mas aponte que o resumo é preferível.
\end{gabarito}

\begin{tarefa}{Tabela de rotas de R3 (CEEI)}
\renewcommand{\arraystretch}{1.0}
\begin{tabular}{|p{3.6cm}|p{2.6cm}|p{2.6cm}|p{5.2cm}|}
\hline
\cabtab{Rede de destino} & \cabtab{Next-hop} & \cabtab{Tipo} & \cabtab{Observação} \\ \hline
\lv \lv \lv \lv \lv \lv \lv
\end{tabular}
\end{tarefa}

\begin{gabarito}
\begin{tabular}{@{}lllp{5.4cm}@{}}
\toprule
\cabtab{Destino} & \cabtab{Next-hop} & \cabtab{Tipo} & \cabtab{Observação} \\
\midrule
10.20.0.0/30   & \cinza{(enlace R1)} & conectada & \\
10.20.0.4/30   & \cinza{(enlace R2)} & conectada & \\
10.20.4.0/25   & 10.20.0.1 & estática & resumo de LCC + LABARC \\
10.20.5.0/27   & 10.20.0.5 & estática & LAEG \\
10.20.5.64/28  & 10.20.0.5 & estática & Lab. Materiais \\
\bottomrule
\end{tabular}

Cinco entradas: duas conectadas, uma resumida para o REENGE e duas específicas para o CAA. Este é o roteador em que a decisão de resumir mais importa, porque é ele que concentra as rotas dos dois prédios.
\end{gabarito}

\subsection{Sumarização}

\begin{pergunta}
Para cada prédio, responda: existe uma \textbf{única} rota que cubra exatamente todas as redes de laboratório daquele prédio, sem incluir nenhum endereço fora delas? Se sim, qual é essa rota? Se não, qual seria o menor bloco que cobre todas elas, e quantos endereços a mais ele incluiria? Justifique com o último octeto em binário.
\espacoresposta{7cm}
\end{pergunta}

\begin{gabarito}
\textbf{REENGE.} LCC (\codigo{.0} a \codigo{.63}) e LABARC (\codigo{.64} a \codigo{.127}). Em binário, \codigo{0 = 00000000} e \codigo{64 = 01000000}: os dois blocos compartilham o primeiro bit do último octeto, e o resto é livre. Resumo exato: \codigo{10.20.4.0/25} (128 endereços, todos pertencentes a um dos dois laboratórios).

\textbf{CAA.} LAEG (\codigo{.0} a \codigo{.31}) e Lab. Materiais (\codigo{.64} a \codigo{.79}). \codigo{0 = 00000000} e \codigo{64 = 01000000} compartilham apenas o primeiro bit, então o menor bloco que cobre os dois é \codigo{10.20.5.0/25} (128 endereços). Em uso: $32 + 16 = 48$. A rota incluiria 80 endereços que não pertencem a nenhum laboratório (o buraco \codigo{.32} a \codigo{.63} e a faixa \codigo{.80} a \codigo{.127}). Não há resumo exato: duas rotas específicas.
\end{gabarito}

\section{Expansão da rede}

Passado algum tempo, dois novos laboratórios entraram em operação: o \termo{LABSPI} (Laboratório de Segurança e Privacidade da Informação) no REENGE e o \termo{LTD} (Laboratório de Técnicas Digitais) no CAA. Os blocos de endereço já estavam reservados desde o planejamento inicial; agora é preciso revisar as tabelas de rotas dos três roteadores.

\begin{center}
\begin{tikzpicture}
  % Roteadores
  \node[roteador, label={[rotulo]right:{R3 (CEEI)}}]  (r3) at (0,2.8) {};
  \node[roteador, label={[rotulo]left:{R1 (REENGE)}}] (r1) at (-4.4,0) {};
  \node[roteador, label={[rotulo]right:{R2 (CAA)}}]   (r2) at (4.4,0) {};

  % Laboratorios do REENGE
  \host{lcc}{(-6.4,-2.6)}{LCC\\10.20.4.0/26}
  \host{labarc}{(-4.4,-2.6)}{LABARC\\10.20.4.64/26}
  \dispositivo[dispositivo,dashed]{\faLaptop}{labspi}{(-2.4,-2.6)}{LABSPI (novo)\\10.20.4.224/28}

  % Laboratorios do CAA
  \host{laeg}{(2.4,-2.6)}{LAEG\\10.20.5.0/27}
  \host{mat}{(4.4,-2.6)}{Lab. Materiais\\10.20.5.64/28}
  \dispositivo[dispositivo,dashed]{\faLaptop}{ltd}{(6.4,-2.6)}{LTD (novo)\\10.20.5.80/28}

  % Enlaces com o CEEI
  \draw[tronco] (r1) -- node[rotuloenlace]{10.20.0.0/30} (r3);
  \draw[tronco] (r2) -- node[rotuloenlace]{10.20.0.4/30} (r3);

  % Enlace de contingencia (fora do escopo deste roteiro)
  \draw[wan] (r1) -- node[rotuloenlace,below]{contingência 10.20.0.8/30} (r2);

  % Enlaces de laboratorio
  \draw[enlace] (r1) -- (lcc);
  \draw[enlace] (r1) -- (labarc);
  \draw[enlace,dashed] (r1) -- (labspi);
  \draw[enlace] (r2) -- (laeg);
  \draw[enlace] (r2) -- (mat);
  \draw[enlace,dashed] (r2) -- (ltd);

  % Dominios
  \begin{scope}[on background layer]
    \node[dominio, fit={(r1)(lcc)(labspi)($(lcc.west)+(-0.5,0)$)($(labspi.east)+(0.6,0)$)($(labarc.south)+(0,-1.1)$)}] (dreenge) {};
    \node[dominio, fit={(r2)(laeg)(ltd)($(laeg.west)+(-0.6,0)$)($(ltd.east)+(0.5,0)$)($(mat.south)+(0,-1.1)$)}] (dcaa) {};
    \node[dominio, fit={(r3)($(r3.north)+(0,0.5)$)($(r3.west)+(-1.4,0)$)($(r3.east)+(2.2,0)$)}] (dceei) {};
  \end{scope}
  \node[dominiotitulo] at (dreenge.north west) {REENGE \quad 10.20.4.0/24};
  \node[dominiotitulo] at (dcaa.north west)    {CAA \quad 10.20.5.0/24};
  \node[dominiotitulo] at (dceei.north west)   {CEEI (agregação)};
\end{tikzpicture}
\end{center}

Os laboratórios tracejados (LABSPI\footnote{Laboratório de Segurança e Privacidade da Informação - DEE} e LTD\footnote{Laboratório de Técnicas Digitais - DEE}) são os que acabaram de entrar em operação. O enlace de contingência entre R1 e R2 continua fora do escopo deste roteiro.

\subsection{Endereçamento}

Novos blocos, já definidos pela equipe de TI:

\begin{center}
\renewcommand{\arraystretch}{1.25}
\begin{tabular}{@{}llll@{}}
\toprule
\cabtab{Prédio} & \cabtab{Laboratório} & \cabtab{Bloco} & \cabtab{Gateway} \\
\midrule
REENGE & LABSPI (Segurança e Privacidade da Informação) & 10.20.4.224/28 & 10.20.4.225 \\
CAA    & LTD (Técnicas Digitais)                         & 10.20.5.80/28  & 10.20.5.81  \\
\bottomrule
\end{tabular}
\end{center}

O LABSPI entra no REENGE e o LTD entra no CAA, com os blocos da tabela acima. Refaça as três tabelas. Reaproveite o que continua valendo do cenário anterior e marque na coluna Observação o que mudou.

\begin{tarefa}{Tabela de rotas de R1 (REENGE), atualizada}
\renewcommand{\arraystretch}{1.0}
\begin{tabular}{|p{3.6cm}|p{2.6cm}|p{2.6cm}|p{5.2cm}|}
\hline
\cabtab{Rede de destino} & \cabtab{Next-hop} & \cabtab{Tipo} & \cabtab{Observação} \\ \hline
\lv \lv \lv \lv \lv \lv \lv \lv
\end{tabular}
\end{tarefa}

\begin{gabarito}
\begin{tabular}{@{}lllp{5.4cm}@{}}
\toprule
\cabtab{Destino} & \cabtab{Next-hop} & \cabtab{Tipo} & \cabtab{Observação} \\
\midrule
10.20.4.0/26   & \cinza{(VLAN LCC)}    & conectada & \\
10.20.4.64/26  & \cinza{(VLAN LABARC)} & conectada & \\
10.20.4.224/28 & \cinza{(VLAN LABSPI)} & conectada & nova \\
10.20.0.0/30   & \cinza{(enlace CEEI)} & conectada & \\
10.20.5.0/27   & 10.20.0.2 & estática & LAEG, sem mudança \\
10.20.5.64/27  & 10.20.0.2 & estática & resumo novo: Lab. Materiais + LTD (substitui o /28) \\
\bottomrule
\end{tabular}
\end{gabarito}

\begin{tarefa}{Tabela de rotas de R2 (CAA), atualizada}
\renewcommand{\arraystretch}{1.0}
\begin{tabular}{|p{3.6cm}|p{2.6cm}|p{2.6cm}|p{5.2cm}|}
\hline
\cabtab{Rede de destino} & \cabtab{Next-hop} & \cabtab{Tipo} & \cabtab{Observação} \\ \hline
\lv \lv \lv \lv \lv \lv \lv \lv
\end{tabular}
\end{tarefa}

\begin{gabarito}
\begin{tabular}{@{}lllp{5.4cm}@{}}
\toprule
\cabtab{Destino} & \cabtab{Next-hop} & \cabtab{Tipo} & \cabtab{Observação} \\
\midrule
10.20.5.0/27   & \cinza{(VLAN LAEG)}      & conectada & \\
10.20.5.64/28  & \cinza{(VLAN Materiais)}  & conectada & \\
10.20.5.80/28  & \cinza{(VLAN LTD)}        & conectada & nova \\
10.20.0.4/30   & \cinza{(enlace CEEI)}     & conectada & \\
10.20.4.0/25   & 10.20.0.6 & estática & resumo LCC + LABARC, sem mudança \\
10.20.4.224/28 & 10.20.0.6 & estática & LABSPI, nova e específica (não cabe no /25) \\
\bottomrule
\end{tabular}
\end{gabarito}

\begin{tarefa}{Tabela de rotas de R3 (CEEI), atualizada}
\renewcommand{\arraystretch}{1.0}
\begin{tabular}{|p{3.6cm}|p{2.6cm}|p{2.6cm}|p{5.2cm}|}
\hline
\cabtab{Rede de destino} & \cabtab{Next-hop} & \cabtab{Tipo} & \cabtab{Observação} \\ \hline
\lv \lv \lv \lv \lv \lv \lv \lv
\end{tabular}
\end{tarefa}

\begin{gabarito}
\begin{tabular}{@{}lllp{5.4cm}@{}}
\toprule
\cabtab{Destino} & \cabtab{Next-hop} & \cabtab{Tipo} & \cabtab{Observação} \\
\midrule
10.20.0.0/30   & \cinza{(enlace R1)} & conectada & \\
10.20.0.4/30   & \cinza{(enlace R2)} & conectada & \\
10.20.4.0/25   & 10.20.0.1 & estática & LCC + LABARC, sem mudança \\
10.20.4.224/28 & 10.20.0.1 & estática & LABSPI, nova e específica \\
10.20.5.0/27   & 10.20.0.5 & estática & LAEG, sem mudança \\
10.20.5.64/27  & 10.20.0.5 & estática & Lab. Materiais + LTD, resumo novo \\
\bottomrule
\end{tabular}

Os dois prédios terminam com duas rotas cada em R3, mas por caminhos opostos: o REENGE tinha um resumo e ganhou uma específica; o CAA tinha duas específicas e passou a ter um resumo mais uma específica. Isso mostra que resumir depende do alinhamento dos blocos, não de "pertencer ao mesmo prédio".
\end{gabarito}

\subsection{Sumarização}

\begin{pergunta}
Agora com os laboratórios novos: para cada prédio, existe uma \textbf{única} rota que cubra exatamente todas as redes de laboratório daquele prédio? Se não, qual seria o menor bloco que cobre todas elas, e quantos endereços a mais ele incluiria? Justifique com o último octeto em binário.
\espacoresposta{7cm}
\end{pergunta}

\begin{gabarito}
\textbf{REENGE.} Acrescenta LABSPI (\codigo{.224} a \codigo{.239}). \codigo{0 = 00000000} e \codigo{224 = 11100000} diferem já no primeiro bit, então o menor bloco que cobre os três é o \codigo{/24} inteiro (256 endereços) para $64 + 64 + 16 = 144$ em uso: 112 endereços a mais. Não há resumo exato para os três. A melhor tabela mantém \codigo{10.20.4.0/25} (resumo exato de LCC + LABARC) e acrescenta \codigo{10.20.4.224/28}.

\textbf{CAA.} Acrescenta LTD (\codigo{.80} a \codigo{.95}). Lab. Materiais e LTD: \codigo{64 = 01000000} e \codigo{80 = 01010000} compartilham os três primeiros bits, então \codigo{10.20.5.64/27} (\codigo{.64} a \codigo{.95}) cobre exatamente os dois. Com a LAEG junto, o menor bloco volta a ser \codigo{10.20.5.0/25} (128 endereços para $32 + 16 + 16 = 64$ em uso: 64 a mais). Melhor tabela: \codigo{10.20.5.0/27} mais \codigo{10.20.5.64/27}.

\textbf{Nota para a discussão.} Em operação real, é comum anunciar o \codigo{/24} inteiro de cada prédio, aceitando que a rota cubra endereços não usados, em troca de uma tabela menor e estável quando novos laboratórios aparecem. O roteiro pede o resumo exato para exercitar o cálculo; se um aluno defender o \codigo{/24} com essa justificativa, a resposta é válida, desde que ele saiba dizer quantos endereços a mais está cobrindo.
\end{gabarito}

% ---------------------------------------------------------------------------
% >>> OPCIONAL, REVISAR ANTES DE LIBERAR: pergunta de isolamento.
% Conceitual, sem exigir implementacao. Descomente o bloco abaixo para incluir.
% ---------------------------------------------------------------------------
% \section{Isolamento}
%
% \begin{pergunta}
% No REENGE, LCC e LABARC ficam atrás do mesmo roteador. Com as tabelas que você
% montou, um computador do LCC consegue alcançar um computador do LABARC? Se a
% coordenação pedisse que as duas redes ficassem isoladas uma da outra (sem
% tráfego entre elas), mantendo as duas com acesso ao CEEI, bastaria mexer nas
% rotas? O que mais seria necessário? Responda em poucas linhas, sem precisar
% indicar comandos.
% \espacoresposta{4cm}
% \end{pergunta}
%
% \begin{gabarito}
% Sim, alcança: as duas redes são conectadas em R1, e o roteador encaminha entre
% elas assim que o roteamento está habilitado. Rotas conectam, não isolam. Para
% impedir o tráfego entre LCC e LABARC mantendo o acesso ao CEEI, seria preciso
% um mecanismo de filtragem (ACL no roteador ou no switch), não uma mudança de
% rota. Gancho com o Módulo 03 (regras no switch OS).
% \end{gabarito}

% ---------------------------------------------------------------------------
% >>> OPCIONAL, REVISAR ANTES DE LIBERAR: rota padrao em R1 e R2.
% ---------------------------------------------------------------------------
% \begin{pergunta}
% Se a restrição de não usar rota padrão fosse retirada, R1 e R2 poderiam
% substituir todas as suas rotas estáticas por uma única rota \codigo{0.0.0.0/0}
% apontando para R3? Que vantagem isso traria? O que continuaria sendo
% necessário em R3?
% \espacoresposta{4cm}
% \end{pergunta}
%
% \begin{gabarito}
% Sim. Vantagem: uma única entrada em cada prédio, que não precisa mudar quando
% o outro prédio ganha laboratórios. Em R3 nada muda: ele continua precisando
% de rotas explícitas (resumidas ou específicas) para cada prédio, porque é ele
% quem decide para qual enlace mandar cada destino. A decisão de resumir não
% desaparece, só se concentra no ponto de agregação.
% \end{gabarito}

% \begin{entrega}
% As seis tabelas preenchidas (três para o cenário atual e três para a expansão) e as respostas das duas perguntas de sumarização, em papel ou digitalizado. Este roteiro não envolve configuração de equipamento: o mesmo cenário será montado no GNS3 no roteiro prático seguinte, e as tabelas daqui servem de referência para as rotas que você vai configurar lá.
% \end{entrega}

\end{document}
